\section{환 준동형과 몫환}
\label{sec:ring-homomorphisms-quotients}

\begin{learninggoal}
환 준동형의 핵과 상을 계산하고, 아이디얼로 몫환을 구성하며, 환의 제1동형정리를 실제 계산에 적용한다.
\end{learninggoal}

\subsection{연산을 보존하는 함수}

\begin{definition}[환 준동형]
환 $R,S$ 사이의 함수 $\varphi:R\to S$가 모든 $a,b\in R$에 대해
\[
\begin{aligned}
 \varphi(a+b)&=\varphi(a)+\varphi(b),\\
 \varphi(ab)&=\varphi(a)\varphi(b),\\
 \varphi(1_R)&=1_S
\end{aligned}
\]
를 만족하면 $\varphi$를 \emph{환 준동형}(ring homomorphism)이라 한다.
\end{definition}

\begin{proposition}
환 준동형 $\varphi:R\to S$에 대해 다음이 성립한다.
\begin{enumerate}[label=(\roman*)]
  \item $\varphi(0)=0$.
  \item $\varphi(-a)=-\varphi(a)$.
  \item $\ker\varphi=\{a\in R:\varphi(a)=0\}$는 $R$의 아이디얼이다.
  \item $\im\varphi=\varphi(R)$는 $S$의 부분환이다.
  \item $u\in R^\times$이면 $\varphi(u)\in S^\times$이고 $\varphi(u^{-1})=\varphi(u)^{-1}$이다.
\end{enumerate}
\end{proposition}

\begin{proof}
(i), (ii)는 덧셈군 준동형의 성질이다. $a,b\in\ker\varphi$이면
\[
  \varphi(a-b)=\varphi(a)-\varphi(b)=0
\]
이고, $r\in R$에 대해
\[
  \varphi(ra)=\varphi(r)\varphi(a)=0
\]
이므로 핵은 아이디얼이다. 상이 부분환이라는 사실은 연산 보존성과 $\varphi(1_R)=1_S$에서 따른다. 마지막으로 $uu^{-1}=1$에 $\varphi$를 적용하면 된다.
\end{proof}

\begin{example}[정수환의 보편성]
임의의 환 $R$에 대해 환 준동형
\[
  \iota_R:\Z\to R,\qquad n\mapsto n\cdot1_R
\]
가 정확히 하나 존재한다. 그 핵은 어떤 $n\ge0$에 대해 $n\Z$ 꼴이다. 이 $n$을 $R$의 \emph{표수}라 하며 $\operatorname{char}R$로 쓴다. 표수는 체확장을 공부할 때 본격적으로 다룬다.
\end{example}

\begin{example}[평가 준동형]
$R\subseteq S$가 환의 확장이고 $\alpha\in S$라 하자. 다항식 평가
\[
  \operatorname{ev}_\alpha:R[X]\to S,
  \qquad f\mapsto f(\alpha)
\]
는 환 준동형이다. 핵은 $\alpha$를 근으로 갖는 $R[X]$의 모든 다항식으로 이루어진다.
\end{example}

\begin{pitfall}
환 준동형의 상은 부분환이지만 일반적으로 아이디얼은 아니다. 예를 들어 포함사상 $\Z\hookrightarrow\Q$의 상 $\Z$는 $\Q$의 아이디얼이 아니다. 실제로 $1\in\Z$이므로 아이디얼이라면 $\Z=\Q$여야 한다.
\end{pitfall}

\subsection{아이디얼로 원소를 0으로 만들기}

아이디얼 $I\trianglelefteq R$를 하나의 ``영으로 취급할 부분''으로 생각하자. $a,b\in R$에 대해
\[
  a\sim b\quad\Longleftrightarrow\quad a-b\in I
\]
로 두면 동치관계가 된다. $a$의 동치류를 $a+I$로 쓴다.

\begin{definition}[몫환]
아이디얼 $I\trianglelefteq R$에 대해 잉여류 집합
\[
  R/I=\{a+I:a\in R\}
\]
에
\[
\begin{aligned}
 (a+I)+(b+I)&=(a+b)+I,\\
 (a+I)(b+I)&=ab+I
\end{aligned}
\]
로 연산을 정의한다. 이렇게 얻은 환을 $R$의 $I$에 의한 \emph{몫환}(quotient ring)이라 한다.
\end{definition}

곱셈이 잘 정의된다는 점이 핵심이다. $a'=a+i$, $b'=b+j$이고 $i,j\in I$라면
\[
  a'b'-ab=aj+bi+ij\in I
\]
이므로 $a'b'+I=ab+I$이다.

\begin{proposition}
사상
\[
  \pi:R\to R/I,\qquad a\mapsto a+I
\]
는 전사 환 준동형이며 $\ker\pi=I$이다.
\end{proposition}

\begin{example}
\begin{enumerate}[label=(\alph*)]
  \item $\Z/(n)=\Z/n\Z$.
  \item $R[X]/(X)$에서는 $X=0$으로 취급되므로 $R[X]/(X)\cong R$이다.
  \item $R[X]/(X-a)\cong R$이며 동형은 $f+(X-a)\mapsto f(a)$로 주어진다.
  \item $\R[X]/(X^2+1)\cong\C$이다. 몫환에서 $X^2=-1$이 되므로 $X$의 잉여류가 허수단위의 역할을 한다.
\end{enumerate}
\end{example}

\subsection{몫환의 보편성}

\begin{theorem}[환의 준동형 기본정리]
환 준동형 $\varphi:R\to S$와 아이디얼 $I\subseteq\ker\varphi$가 주어졌다고 하자. 그러면 유일한 환 준동형
\[
  \overline\varphi:R/I\to S
\]
가 존재하여
\[
  \overline\varphi(a+I)=\varphi(a)
\]
를 만족한다. 즉 $\varphi=\overline\varphi\circ\pi$이다.
\end{theorem}

\begin{proofidea}
$I$의 원소를 모두 $0$으로 보내는 준동형은 같은 잉여류에 속한 원소들을 구별하지 못한다. 따라서 잉여류 $a+I$를 $\varphi(a)$로 보내는 것이 자연스럽다. 핵심은 대표원 선택과 무관함을 확인하는 것이다.
\end{proofidea}

\begin{proof}
$a+I=b+I$이면 $a-b\in I\subseteq\ker\varphi$이므로
\[
  \varphi(a)-\varphi(b)=\varphi(a-b)=0.
\]
따라서 $\overline\varphi(a+I)=\varphi(a)$는 잘 정의된다. 덧셈, 곱셈, 항등원을 보존한다는 것은 $\varphi$의 성질에서 바로 따른다. $\pi$가 전사이므로 이런 사상은 유일하다.
\end{proof}

\begin{corollary}[제1동형정리]
환 준동형 $\varphi:R\to S$에 대해 자연스러운 환 동형
\[
  R/\ker\varphi\cong\im\varphi
\]
가 존재한다.
\end{corollary}

\begin{example}[평가와 인수정리의 구조적 형태]
체 $K$와 $a\in K$에 대해
\[
  \operatorname{ev}_a:K[X]\to K,\qquad f\mapsto f(a)
\]
는 전사이다. 나눗셈 알고리즘에 의해 $f(a)=0$일 필요충분조건은 $(X-a)\mid f$이므로
\[
  \ker(\operatorname{ev}_a)=(X-a).
\]
따라서 제1동형정리에 의해
\[
  K[X]/(X-a)\cong K.
\]
\end{example}

\begin{example}[복소수의 대수적 구성]
사상
\[
  \Phi:\R[X]\to\C,\qquad f\mapsto f(i)
\]
를 생각하자. 모든 복소수 $a+bi$는 $a+bX$의 상이므로 $\Phi$는 전사이다. $X^2+1$은 핵에 속하고, 다항식 나눗셈으로 핵이 정확히 $(X^2+1)$임을 알 수 있다. 따라서
\[
  \R[X]/(X^2+1)\cong\C.
\]
이 예는 ``방정식 $X^2+1=0$의 해를 새로 붙인다''는 체확장의 기본 아이디어를 보여 준다.
\end{example}

\subsection{환의 동형정리}

\begin{theorem}[제2동형정리]
$R$의 부분환 $S$와 아이디얼 $I\trianglelefteq R$에 대해 $S+I$는 $R$의 부분환이고 $S\cap I$는 $S$의 아이디얼이며
\[
  S/(S\cap I)\cong(S+I)/I
\]
이다.
\end{theorem}

\begin{proof}
합성
\[
  S\hookrightarrow S+I\twoheadrightarrow(S+I)/I
\]
은 전사 환 준동형이고 핵은 $S\cap I$이다. 제1동형정리를 적용한다.
\end{proof}

\begin{theorem}[제3동형정리]
$I\subseteq J$가 $R$의 아이디얼이면 $J/I$는 $R/I$의 아이디얼이고
\[
  (R/I)/(J/I)\cong R/J
\]
이다.
\end{theorem}

\begin{checkpoint}
다음 동형을 적절한 준동형의 핵을 계산하여 증명하라.
\[
  \Z[X]/(2,X)\cong\F_2,
  \qquad
  \Q[X]/(X^2-2)\cong\Q[\sqrt2].
\]
\end{checkpoint}
